\subsection{普遍可测函数}

定义 2. --- T 到拓扑空间 F 中的映射 f 称为普遍可测的，如果它对 T 上每个正测度 \(\mu\) 都是 \(\mu\) -可测的。

T 的子集中，其特征函数为普遍可测者，称为普遍可测集。它们构成 T 上的一个 σ-集环（tribe）（Ch. IV, §5, No.~4, 定理 2 的推论 2），该 σ-集环包含 Borel 集（同处引证，推论 3），并且当 T 可度量化时还包含 Souslin 集（Ch. IV, §5, No.~1, 命题 3 的推论 2）。T 到可度量化且可分的拓扑空间 F 中的映射 f 为普遍可测的，必须且只须 F 中每个闭球在 f 下的逆像是 T 的普遍可测子集（Ch. IV, §5, No.~5, 定理 4）。

命题 6. --- T 到拓扑空间 F 中的映射 f 为普遍可测的，必须且只须 f 对 T 上每个具有紧支集的正测度都可测。

该条件显然是必要的；另一方面它也是充分的，因为每个正测度 \(\mu\) 都是一族具有紧支集的测度之和（ \(\S2\) ，No.~3，命题 4）：于是结论由 \(\S2\) ，No.~2 的命题 2 得出。

命题 7. --- 设 \(\mu\) 是 T 上的一个正测度，\(f\) 是 T 到拓扑空间 F 中的一个 \(\mu\) -可测映射。则存在 T 到 F 中的一个普遍可测映射 \(f'\) ，使得 \(f = f'\) 局部 \(\mu\) -几乎处处成立。

设 \(\mathfrak{K}\) 是 T 中使 \(f\) 在 K 上的限制为连续的紧集组成的集；由于 \(\mathfrak{K}\) 是 \(\mu\) -稠密的（Ch. IV, §5, No.~10, 命题 15），存在一个局部可数族 \((\mathrm{K}_i)_{i\in I}\) （由 \(\mathfrak{K}\) 中两两不相交的元素组成），使得集 \(\mathrm{N} = \mathrm{T} - \bigcup_{i\in I}\mathrm{K}_i\) 是局部 \(\mu\) -可忽略的（Ch. IV, §5, No.~9,

命题 14）。设 \(x\) 是 \(\mathbf{F}\) 的一个元素；置

\[
f ^ {\prime} (t) = f (t) \quad \text { if } t \in \bigcup_ {i \in I} K _ {i},
\]

\[
f ^ {\prime} (t) = x \quad \text { if } t \in \mathbb {N}.
\]

函数 \(f\) 与 \(f'\) 局部 \(\mu\) -几乎处处相等。另一方面，\(\mathrm{N} \cap \mathrm{K}\) 是 \(\mathrm{K}\) 的一个 Borel 子集，其中 \(\mathrm{K}\) 是 \(\mathrm{T}\) 中任意的紧集，因为族 \((\mathrm{K}_i)\) 是局部可数的。由此得出 \(\mathrm{N}\) 是普遍可测集，且 \(f'\) 是普遍可测函数（Ch. IV, §5, No.~10, 命题 16）。

